\documentclass{amsart}
% For dealing with references we use the comment environment
\usepackage{verbatim}
\newenvironment{reference}{\comment}{\endcomment}
%\newenvironment{reference}{}{}
\newenvironment{slogan}{\comment}{\endcomment}
\newenvironment{history}{\comment}{\endcomment}

% For commutative diagrams we use Xy-pic
\usepackage[all]{xy}

% We use 2cell for 2-commutative diagrams.
\xyoption{2cell}
\UseAllTwocells

% We use multicol for the list of chapters between chapters
\usepackage{multicol}

% This is generally recommended for better output
\usepackage{lmodern}
\usepackage[T1]{fontenc}

% For cross-file-references

% Package for hypertext links:
\usepackage{hyperref}

% For any local file, say "hello.tex" you want to link to please

% Theorem environments.
%
\theoremstyle{plain}
\newtheorem{theorem}[subsection]{Theorem}
\newtheorem{proposition}[subsection]{Proposition}
\newtheorem{lemma}[subsection]{Lemma}

\theoremstyle{definition}
\newtheorem{definition}[subsection]{Definition}
\newtheorem{example}[subsection]{Example}
\newtheorem{exercise}[subsection]{Exercise}
\newtheorem{situation}[subsection]{Situation}

\theoremstyle{remark}
\newtheorem{remark}[subsection]{Remark}
\newtheorem{remarks}[subsection]{Remarks}

\numberwithin{equation}{subsection}

% Macros
%
\def\lim{\mathop{\mathrm{lim}}\nolimits}
\def\colim{\mathop{\mathrm{colim}}\nolimits}
\def\Spec{\mathop{\mathrm{Spec}}}
\def\Hom{\mathop{\mathrm{Hom}}\nolimits}
\def\Ext{\mathop{\mathrm{Ext}}\nolimits}
\def\SheafHom{\mathop{\mathcal{H}\!\mathit{om}}\nolimits}
\def\SheafExt{\mathop{\mathcal{E}\!\mathit{xt}}\nolimits}
\def\Sch{\mathit{Sch}}
\def\Mor{\mathop{\mathrm{Mor}}\nolimits}
\def\Ob{\mathop{\mathrm{Ob}}\nolimits}
\def\Sh{\mathop{\mathit{Sh}}\nolimits}
\def\NL{\mathop{N\!L}\nolimits}
\def\CH{\mathop{\mathrm{CH}}\nolimits}
\def\proetale{{pro\text{-}\acute{e}tale}}
\def\etale{{\acute{e}tale}}
\def\QCoh{\mathit{QCoh}}
\def\Ker{\mathop{\mathrm{Ker}}}
\def\Im{\mathop{\mathrm{Im}}}
\def\Coker{\mathop{\mathrm{Coker}}}
\def\Coim{\mathop{\mathrm{Coim}}}

% Boxtimes
%
\DeclareMathSymbol{\boxtimes}{\mathbin}{AMSa}{"02}

%
% Macros for moduli stacks/spaces
%
\def\QCohstack{\mathcal{QC}\!\mathit{oh}}
\def\Cohstack{\mathcal{C}\!\mathit{oh}}
\def\Spacesstack{\mathcal{S}\!\mathit{paces}}
\def\Quotfunctor{\mathrm{Quot}}
\def\Hilbfunctor{\mathrm{Hilb}}
\def\Curvesstack{\mathcal{C}\!\mathit{urves}}
\def\Polarizedstack{\mathcal{P}\!\mathit{olarized}}
\def\Complexesstack{\mathcal{C}\!\mathit{omplexes}}
% \Pic is the operator that assigns to X its picard group, usage \Pic(X)
% \Picardstack_{X/B} denotes the Picard stack of X over B
% \Picardfunctor_{X/B} denotes the Picard functor of X over B
\def\Pic{\mathop{\mathrm{Pic}}\nolimits}
\def\Picardstack{\mathcal{P}\!\mathit{ic}}
\def\Picardfunctor{\mathrm{Pic}}
\def\Deformationcategory{\mathcal{D}\!\mathit{ef}}

\setlength{\emergencystretch}{3em}
\begin{document}
\title[Stacks Project, Tag 09Z0]{481 --- Étale covers of quasi-compact quasi-separated schemes admit finite finitely presented surjective refinements that factor locally through the cover}
\maketitle
\noindent Copyright (C) 2005--2025 Johan de Jong. Stacks Project authors. Source: \href{https://stacks.math.columbia.edu/tag/09Z0}{Tag 09Z0}.

\noindent Editorial difficulty note (not author text). Difficulty H2: Induction on geometric fibre cardinalities must construct an integral cover despite the boundary of an étale open immersion. The proof separates the diagonal component, extends a smaller-degree refinement across the boundary, and finally descends its local factorization data to a finite finitely presented approximation..

\noindent Editorial provenance note (not author text). History: excerpt from revision \texttt{a0444\allowbreak 6e57e\allowbreak c1fbc\allowbreak 252a8\allowbreak 71afc\allowbreak ec775\allowbreak 2fb28\allowbreak 07b14}, more-morphisms.tex, lines 12978--13070. Reference presentation was adapted; original-author context and core support blocks were retained or added. The mathematical wording is unchanged. Licensed under GNU FDL 1.2 or later; no invariant sections or cover texts. See ../licenses/COPYING. Context blocks reproduce author definitions and supporting results; they are not additional counted problems.

\begin{lemma}
\label{lemma-there-is-a-scheme-integral-over}
Let $U \to X$ be a surjective \'etale morphism of schemes. Assume $X$
is quasi-compact and quasi-separated. Then there exists a surjective
integral morphism $Y \to X$, such that for
every $y \in Y$ there is an open neighbourhood $V \subset Y$
such that $V \to X$ factors through $U$. In fact, we may assume
$Y \to X$ is finite and of finite presentation.
\end{lemma}

\begin{proof}
Since $X$ is quasi-compact, there exist finitely many affine opens
$U_i \subset U$ such that $U' = \coprod U_i \to X$ is surjective.
After replacing $U$ by $U'$, we see that we may assume $U$ is affine.
In particular $U \to X$ is separated
(Schemes, Lemma \href{https://stacks.math.columbia.edu/tag/01KN}{lemma affine separated (Tag 01KN)}).
Then there exists an integer $d$ bounding the degree of the geometric
fibres of $U \to X$ (see Morphisms, Lemma
\href{https://stacks.math.columbia.edu/tag/03JA}{lemma locally quasi finite qc source universally bounded (Tag 03JA)}).
We will prove the lemma by induction on $d$ for all quasi-compact
and separated schemes $U$ mapping surjective and \'etale onto $X$.
If $d = 1$, then $U = X$ and the result holds with $Y = U$.
Assume $d > 1$.

\medskip\noindent
We apply Lemma \href{https://stacks.math.columbia.edu/tag/02LR}{lemma quasi finite separated quasi affine (Tag 02LR)}
and we obtain a factorization
$$
\xymatrix{
U \ar[rr]_j \ar[rd] & & Y \ar[ld]^\pi \\
& X
}
$$
with $\pi$ integral and $j$ a quasi-compact open immersion. We may and do
assume that $j(U)$ is scheme theoretically dense in $Y$. Note that
$$
U \times_X Y = U \amalg W
$$
where the first summand is the image of $U \to U \times_X Y$
(which is closed by
Schemes, Lemma \href{https://stacks.math.columbia.edu/tag/01KS}{lemma semi diagonal (Tag 01KS)}
and open because it is \'etale as a morphism between schemes \'etale over $Y$)
and the second summand is the (open and closed) complement.
The image $V \subset Y$ of $W$ is an open subscheme containing
$Y \setminus U$.

\medskip\noindent
The \'etale morphism $W \to Y$ has geometric fibres of cardinality $< d$.
Namely, this is clear for geometric points of $U \subset Y$ by inspection.
Since $U \subset Y$ is dense, it holds for all geometric points of $Y$
for example by Lemma
\href{https://stacks.math.columbia.edu/tag/07RY}{lemma stratify flat fp lqf universally bounded (Tag 07RY)}
(the degree of the fibres of a quasi-compact separated \'etale morphism
does not go up under specialization). Thus we may apply the induction
hypothesis to $W \to V$ and find a surjective integral morphism
$Z \to V$ with $Z$ a scheme, which Zariski locally factors through $W$.
Choose a factorization $Z \to Z' \to Y$ with $Z' \to Y$ integral and
$Z \to Z'$ open immersion
(Lemma \href{https://stacks.math.columbia.edu/tag/02LR}{lemma quasi finite separated quasi affine (Tag 02LR)}).
After replacing $Z'$ by the scheme theoretic closure of $Z$ in $Z'$
we may assume that $Z$ is scheme theoretically dense in $Z'$.
After doing this we have $Z' \times_Y V = Z$. Finally,
let $T \subset Y$ be the induced reduced closed subscheme structure
on $Y \setminus V$. Consider the morphism
$$
Z' \amalg T \longrightarrow X
$$
This is a surjective integral morphism by construction.
Since $T \subset U$ it is clear that the morphism $T \to X$
factors through $U$. On the other hand, let $z \in Z'$
be a point. If $z \not \in Z$, then $z$ maps to a point of
$Y \setminus V \subset U$ and we find a neighbourhood of $z$
on which the morphism factors through $U$.
If $z \in Z$, then we have a neighbourhood $\Omega \subset Z$
which factors through $W \subset U \times_X Y$ and hence through $U$.
This proves existence.

\medskip\noindent
Assume we have found $Y \to X$ integral and surjective which Zariski
locally factors through $U$. Choose a finite affine open covering
$Y = \bigcup V_j$ such that $V_j \to X$ factors through $U$. We can
write $Y = \lim Y_i$ with $Y_i \to X$ finite and of finite
presentation, see Limits, Lemma
\href{https://stacks.math.columbia.edu/tag/09YZ}{lemma integral limit finite and finite presentation (Tag 09YZ)}.
For large enough $i$ we can find affine opens $V_{i, j} \subset Y_i$
whose inverse image in $Y$ recovers $V_j$, see
Limits, Lemma \href{https://stacks.math.columbia.edu/tag/01Z4}{lemma descend opens (Tag 01Z4)}.
For even larger $i$ the morphisms $V_j \to U$ over $X$ come
from morphisms $V_{i, j} \to U$ over $X$, see
Limits, Proposition
\href{https://stacks.math.columbia.edu/tag/01ZC}{proposition characterize locally finite presentation (Tag 01ZC)}.
This finishes the proof.
\end{proof}
\end{document}
